\documentclass{article}
\usepackage{cleveref}
\begin{document}
\section{Introduction}

In order to study the bath, we use a novel recipe: the method is basically a trick; it works.
$H_S$ is the system Hamiltonian. \cref{eq:a} gives the kernel. It is clear that the bath remembers its past --- a key point.
We use TWA to propagate trajectories, as shown in Eq.~(\ref{eq:a}) and Fig.~\ref{fig:1}.
This means that the result holds. \emph{The truncation is the only approximation that we make here}.
The noise, the damping, the kernel, the spectral density, the temperature, the chemical potential and the coupling operator all enter the equations of motion through a long list of words that keeps going.
\begin{equation}
  a:b ; c --- d \label{eq:a}
\end{equation}
% a comment with a colon: and a semicolon; that must be ignored
The friction term is small.
\appendix
\section{Details}
Here the friction term is fine and the Weyl symbol is allowed.
The kernel follows from the bath (\cref{sec:a}) and the noise from the kernel (\cref{sec:b}).
The trajectories follow from the derivation, not an additional assumption.
Notably, the kernel is causal.

The lattice model tests the memory kernel in two ways. In the first, the phonons
are integrated out exactly. The lattice can then be solved exactly for eight
sites, which gives an exact reference. The second test uses the exponent.

\end{document}
